\section{Related Work}

\paragraph{Input-blind and constant baselines.}
\citet{balepur2024artifacts} define a majority-class baseline for letter MCQA and recommend stronger-than-chance references; their object is dataset cheatability under a black-box generative setup, whereas ours is a per-arm gate for likelihood-scored constructs. The distinction is operational: their invalid-output analysis imputes either 0.0 or 0.25 on MMLU, while the construct's best-constant letter value is 0.2689. \citet{zheng2025cheating} use the term ``null model'' for crafted constant responses that adversarially exploit an LLM judge. Their null is an attack optimized against a judge; ours is a reference derived from benchmark statistics and requires no attacker. The metalinguistic preprint that motivates our opening reports the above-chance/below-majority pattern descriptively on one new benchmark, but does not turn it into a pre-comparison protocol \citep{arcon2026metalinguistic}.

\paragraph{Selection bias and MC interfaces.}
Letter preference is not a new observation. \citet{zheng2024selectors} identify selection and token bias and propose PriDe, which estimates an option-ID prior from permuted contents. PriDe fixes the predictor; null calibration fixes the reference line. Even after prediction debiasing, reporting only against chance can leave the benchmark's gold-marginal floor unreported. We therefore do not reject permutation controls on cost grounds: they are a useful and potentially inexpensive diagnostic. They answer whether predictions should be corrected, while our arm-independent floor answers whether the resulting interface score is comparable across arms.

OLMES standardizes both multiple-choice formulation (MCF) and cloze formulation (CF), then uses the better score per task and model \citep{gu2025olmes}. This is not a size-keyed rule. Its reference discussion is nevertheless framed around a random baseline rather than a measured label-marginal floor, so the max rule can select an interface without testing whether that interface clears an input-blind reference. Our ARC-Easy result shows a case where CF rescues an at-floor letter read-out; MMLU shows that an interface swap need not do so. \citet{cho2026choices} independently study symbols, cloze, and hybrid formats and propose a new question-sensitive score. We do not claim the interface contrast; our method instead gates whichever construct is reported.

\paragraph{Length-based scoring.}
Length-normalized likelihood can over-correct and its token-based form is tokenizer-dependent \citep{oostermeijer2026length}. We therefore do not claim generic length sensitivity or tokenizer dependence of the scoring rule. We show that the \emph{induced input-blind floor} inherits these choices: its numerical value changes with tie convention, character versus continuation-token length, and tokenizer. OLMES's observation that tokenizer dependence may be irrelevant when ranking options for one fixed model and tokenizer is compatible with our result; our scope is cross-model comparison of the reference itself.

\paragraph{Nulls as interpretive controls.}
Control tasks and selectivity can reverse conclusions about which representation layer is informative \citep{hewitt2019probes}. That is the closest structural precedent: their null randomizes supervision for a probe, while ours removes input dependence from an MC construct. \citet{ding2021grounding} similarly motivate sensitivity and specificity tests for representation similarity. Finally, \citet{bean2025measuring} provide the construct-validity vocabulary and a broad benchmark checklist. We position null calibration as the missing operational item, not as the invention of construct validity.
